% Wersja źródłowa: audyt II, jednostka 08
\setcounter{section}{7}
\section{Pary dualne i topologie słabe}\label{sec:topologie-slabe}

\subsection{Dlaczego potrzebujemy kolejnej topologii?}

Dotąd topologia przestrzeni liniowo-topologicznej $X$ była dana z góry,
a przestrzeń dualna topologiczna $X'$ składała się z tych
funkcjonałów liniowych, które są względem niej ciągłe.
Można postąpić odwrotnie: wybrać rodzinę funkcjonałów i skonstruować
\emph{najsłabszą} topologię, w której są one ciągłe.
Z tego pomysłu powstają topologie słabe, kluczowe w badaniu zbieżności
funkcji, operatorów i dystrybucji.

W rozdziale $\K\in\{\R,\C\}$. Wszystkie przestrzenie liniowe
są nad $\K$. Przez $E^*$ oznaczamy przestrzeń dualną algebraiczną,
a przez $E'$ przestrzeń dualną topologiczną, gdy $E$ jest TVS.

\begin{definicja}[Para dualna]\label{def:para-dualna}
Parą dualną nazywamy dwie przestrzenie liniowe $E,F$ wraz
z dwuliniowym parowaniem $\langle\cdot,\cdot\rangle:E\times F\to\K$
takim, że
\begin{enumerate}
\item dla każdego $x\in E\setminus\{0\}$ istnieje $f\in F$
      spełniający $\langle x,f\rangle\ne0$;
\item dla każdego $f\in F\setminus\{0\}$ istnieje $x\in E$
      spełniający $\langle x,f\rangle\ne0$.
\end{enumerate}
Mówimy, że parowanie jest \emph{niezdegenerowane} po obu stronach.
Nie zakładamy istnienia topologii na $E$ ani $F$.
\end{definicja}

\begin{przyklad}[Para kanoniczna]
Jeżeli $E$ jest lokalnie wypukłą TVS Hausdorffa, to
$(E,E')$, z parowaniem $\langle x,f\rangle=f(x)$,
jest parą dualną. Pierwsza niezdegenerowaność wynika
z konsekwencji twierdzenia Hahna--Banacha (\hyperref[sec:hahn-banach]{rozdział 6}),
a druga z faktu, że niezerowy funkcjonał ma niezerową wartość
na pewnym wektorze.
\end{przyklad}

\begin{przyklad}[Para ciągów]
Niech $E=c_0$, a $F=\ell^1$. Parowanie
$\langle x,a\rangle=\sum_{n=1}^{\infty}x_na_n$
jest poprawnie określone, gdyż
$\sum_n|x_na_n|\le\|x\|_\infty\|a\|_1$.
Jest niezdegenerowane: wektory współrzędnych $e_n$ należą
zarówno do $c_0$, jak i do $\ell^1$, i wykrywają odpowiednie
współrzędne.
\end{przyklad}

\subsection{Topologie słabe jako topologie początkowe}

\begin{definicja}[Topologia słaba pary dualnej]\label{def:slaba-para}
Niech $(E,F,\langle\cdot,\cdot\rangle)$ będzie parą dualną.
Topologią $\sigma(E,F)$ nazywamy najsłabszą topologię na $E$,
w której każdy funkcjonał $x\mapsto\langle x,f\rangle$,
$f\in F$, jest ciągły. Analogicznie definiujemy
$\sigma(F,E)$, wykorzystując funkcjonały
$f\mapsto\langle x,f\rangle$ dla $x\in E$.
\end{definicja}

\begin{twierdzenie}[Seminormy i baza otoczeń]\label{thm:slaba-baza}
Niech $(E,F)$ będzie parą dualną. Wtedy $\sigma(E,F)$
jest lokalnie wypukłą topologią Hausdorffa, generowaną przez
seminormy $p_f(x)=|\langle x,f\rangle|$, $f\in F$.
Bazę otoczeń zera stanowią zbiory
\[
 U(f_1,\ldots,f_m;\varepsilon)
 =\{x\in E:|\langle x,f_j\rangle|<\varepsilon,
 \ j=1,\ldots,m\},
\]
gdzie $m\ge1$, $f_1,\ldots,f_m\in F$ i $\varepsilon>0$.
W szczególności sieć $(x_i)_{i\in I}$ zbiega do $x\in E$
w $\sigma(E,F)$ wtedy i tylko wtedy, gdy
$\langle x_i,f\rangle\to\langle x,f\rangle$ dla każdego $f\in F$.
Analogiczne fakty zachodzą dla $\sigma(F,E)$.
\end{twierdzenie}
\begin{proof}
Z definicji topologii początkowej względem rodziny
odwzorowań do $\K$ otrzymujemy zbiory opisane skończenie
wieloma nierównościami. Każda funkcja $p_f$ jest seminormą;
przecięcia jej podpoziomic są absolutnie wypukłe, więc
otrzymana topologia jest lokalnie wypukła. Jeżeli $x\ne0$,
niezdegenerowaność parowania daje $f$ z $p_f(x)>0$.
Można zatem oddzielić $x$ od zera otoczeniem zera,
co dowodzi własności Hausdorffa. Kryterium zbieżności
jest własnością definiującą topologię początkową.
\end{proof}

\begin{stwierdzenie}[Słaba ograniczoność]\label{prop:slaba-ograniczonosc}
Niech $(E,F)$ będzie parą dualną, a $B\subseteq E$.
Zbiór $B$ jest ograniczony w TVS $(E,\sigma(E,F))$
wtedy i tylko wtedy, gdy dla każdego $f\in F$
zachodzi $\sup_{x\in B}|\langle x,f\rangle|<\infty$.
\end{stwierdzenie}
\begin{proof}
Jeżeli $B$ jest ograniczony, jego obraz przez każdy ciągły
funkcjonał jest ograniczonym zbiorem w $\K$.
Odwrotnie, niech $U=U(f_1,\ldots,f_m;\varepsilon)$
będzie podstawowym otoczeniem zera. Z założeń wynika, że
$M=\max_{1\le j\le m}\sup_{x\in B}|\langle x,f_j\rangle|$
jest skończone. Dla każdego $t>\max\{1,M/\varepsilon\}$
mamy $B\subseteq tU$, co jest definicją ograniczoności.
\end{proof}

\begin{twierdzenie}[Przestrzeń dualna topologii słabej]\label{thm:slaby-dual}
Niech $(E,F)$ będzie parą dualną. Odwzorowanie
$F\to(E,\sigma(E,F))'$, przypisujące $f$ funkcjonał
$x\mapsto\langle x,f\rangle$, jest izomorfizmem
algebraicznym. W szczególności topologia słaba nie wytwarza
nowych ciągłych funkcjonałów liniowych poza elementami $F$.
\end{twierdzenie}
\begin{proof}
Niezdegenerowaność zapewnia różnowartościowość, a
ciągłość funkcjonałów $f\in F$ wynika z definicji.
Niech teraz $g:E\to\K$ będzie liniowy i ciągły w
$\sigma(E,F)$. Istnieją $f_1,\ldots,f_m\in F$
i $\varepsilon>0$ takie, że
$|g(x)|<1$ dla $x\in U(f_1,\ldots,f_m;\varepsilon)$.
Połóżmy $T(x)=(\langle x,f_1\rangle,\ldots,
\langle x,f_m\rangle)\in\K^m$. Jeżeli $x\in\ker T$,
to $tx\in U$ dla każdego $t\in\K$; nierówność
$|t g(x)|<1$ dla wszystkich $t$ oznacza $g(x)=0$.
Zatem $g$ faktoryzuje się przez liniową mapę
$\widetilde g:T(E)\to\K$. Przedłużamy ją algebraicznie
z podprzestrzeni $T(E)$ na $\K^m$. Każdy funkcjonał
liniowy na $\K^m$ jest kombinacją współrzędnych,
więc $g(x)=\sum_{j=1}^m a_j\langle x,f_j\rangle
=\langle x,\sum_{j=1}^m a_j f_j\rangle$.
To kończy dowód surjektywności.
\end{proof}

\begin{wniosek}[Topologia wyjściowa i topologia słaba]\label{cor:slaba-wyjsciowa}
Jeżeli $E$ jest lokalnie wypukłą TVS Hausdorffa z topologią
$\tau$ i przestrzenią dualną topologiczną $E'$, to
$\sigma(E,E')\subseteq\tau$ (inkluzja topologii),
a przestrzenie dualne topologiczne $(E,\sigma(E,E'))'$ i
$(E,\tau)'$ są algebraicznie identyczne.
\end{wniosek}
\begin{proof}
Wszystkie elementy $E'$ są ciągłe względem $\tau$,
zatem minimalność topologii początkowej daje inkluzję.
Druga teza wynika z Twierdzenia~\ref{thm:slaby-dual}.
\end{proof}

\subsection{Słabe domknięcia i zjawiska nieskończenie wymiarowe}

\begin{twierdzenie}[Domknięcie podprzestrzeni przez anihilator]
\label{thm:anihilator-domkniecie}
Niech $E$ będzie lokalnie wypukłą TVS Hausdorffa,
$M\subseteq E$ podprzestrzenią liniową, a
$M^\perp=\{f\in E': f(m)=0\ \text{dla wszystkich }m\in M\}$.
Wtedy
\[
\overline M^{\,\tau}=\overline M^{\,\sigma(E,E')}
=\{x\in E:f(x)=0\ \text{dla wszystkich }f\in M^\perp\}.
\]
\end{twierdzenie}
\begin{proof}
Każdy funkcjonał z $M^\perp$ znika na obu domknięciach,
co daje jedną inkluzję. Jeżeli $x_0\notin\overline M^{\,\tau}$,
to w \hyperref[sec:hahn-banach]{rozdziale 6} oddzieliliśmy punkt od domkniętej
podprzestrzeni za pomocą $f\in E'$ spełniającego
$f|_M=0$ i $f(x_0)=1$. Tak więc $x_0$ nie leży w
zbiorze po prawej stronie. Ponieważ domknięcie słabe
zawiera domknięcie w $\tau$, otrzymujemy wszystkie równości.
\end{proof}

\begin{stwierdzenie}[Otoczenia słabe w nieskończonym wymiarze]
\label{prop:slabe-otoczenia}
Niech $(E,F)$ będzie parą dualną, przy czym
$\dim E=\infty$. Wtedy każde otoczenie zera w
$\sigma(E,F)$ zawiera niezerową podprzestrzeń liniową,
a więc nie jest ograniczone w tej topologii.
W konsekwencji $(E,\sigma(E,F))$ nie jest normowalna.
\end{stwierdzenie}
\begin{proof}
Dane otoczenie zawiera $U(f_1,\ldots,f_m;\varepsilon)$,
a ono zawiera $M=\bigcap_{j=1}^m\ker f_j$.
Podprzestrzeń $M$ ma kowymiar nie większy niż $m$,
więc jest niezerowa. Jeżeli $v\in M\setminus\{0\}$,
to $\K v\subseteq U$, lecz zbiór $\K v$ nie jest
słabo ograniczony: wybierając $f\in F$ z
$\langle v,f\rangle\ne0$, otrzymujemy nieograniczone
wartości $|\langle tv,f\rangle|$.
Przestrzeń normowalna ma ograniczone otoczenie zera,
co przeczy powyższemu faktowi.
\end{proof}

\begin{przyklad}[Słaba zbieżność bez zbieżności w normie]
Niech $E=\ell^2$ nad $\R$ albo $\C$ z topologią normową,
a $e_n$ oznacza $n$-ty wektor jednostkowy. Wtedy
$e_n\to0$ słabo, ponieważ każdy funkcjonał ciągły na
$\ell^2$ ma postać $x\mapsto\sum_k x_k a_k$
z pewnym $a\in\ell^2$ (w przypadku zespolonym
odpowiednio wybieramy współrzędne liniowe), a $a_n\to0$.
Jednocześnie $\|e_n\|_2=1$ dla każdego $n$.
\end{przyklad}

\begin{twierdzenie}[Słabe domknięcie sfery jednostkowej]
\label{thm:sfera-slaba}
Niech $E$ będzie nieskończenie wymiarową przestrzenią
unormowaną nad $\K$, $S=\{x\in E:\|x\|=1\}$,
$B=\{x\in E:\|x\|\le1\}$. Wtedy
$\overline S^{\,\sigma(E,E')}=B$.
\end{twierdzenie}
\begin{proof}
Kula $B$ jest słabo domknięta. Istotnie, jeżeli
$\|x_0\|>1$, twierdzenie Hahna--Banacha daje
$f\in E'$ z $\|f\|=1$ i $f(x_0)=\|x_0\|$,
a zbiór $\{x:\operatorname{Re}f(x)>1\}$ jest słabo
otwarty, zawiera $x_0$ i nie przecina $B$.
Stąd $\overline S^{\,\sigma(E,E')}\subseteq B$.

Niech $\|x\|<1$ i $U$ będzie podstawowym słabym
otoczeniem $x$, określonym przez skończenie wiele
funkcjonałów $f_1,\ldots,f_m$. Z nieskończonego wymiaru
wynika istnienie $v\ne0$ we wspólnym jądrze tych
funkcjonałów. Funkcja $t\mapsto\|x+tv\|$, $t\ge0$,
jest ciągła, w zerze mniejsza od $1$, a dla $t\to\infty$
dąży do nieskończoności. Istnieje więc $t_0>0$
z $\|x+t_0v\|=1$. Punkt $x+t_0v$ należy do
$S\cap U$, gdyż funkcjonały definiujące $U$ mają
na nim te same wartości co na $x$. To dowodzi
$B\subseteq\overline S^{\,\sigma(E,E')}$.
\end{proof}

\begin{uwaga}[Co stało się z intuicją geometryczną?]
W przestrzeni skończenie wymiarowej sfera jest domknięta.
W nieskończenie wymiarowej przestrzeni unormowanej topologia
słaba nie rozróżnia, za pomocą skończenie wielu funkcjonałów,
punktów sfery od punktów wewnątrz kuli. To jednak nie znaczy,
że kula jest słabo zwarta: zwartość omówimy osobno.
\end{uwaga}

\subsection{Topologia słaba-* na przestrzeni dualnej}

\begin{definicja}[Topologia słaba-*]\label{def:slaba-gwiazdka}
Niech $E$ będzie lokalnie wypukłą TVS Hausdorffa.
Na przestrzeni dualnej topologicznej $E'$ topologią
\emph{słabą-*} nazywamy $\sigma(E',E)$, tj. topologię
punktowej zbieżności funkcjonałów na $E$. Dla sieci
$(f_i)_{i\in I}\subseteq E'$ zapis
$f_i\overset{*}{\rightharpoonup}f$ oznacza, że
$f_i(x)\to f(x)$ dla każdego $x\in E$.
\end{definicja}

\begin{stwierdzenie}[Słaba i słaba-* topologia na przestrzeni dualnej]
\label{prop:slaba-weakstar}
Niech $E$ będzie przestrzenią Banacha. Na przestrzeni
Banacha $E'$ rozważmy topologię słabą
$\sigma(E',E'')$ i słabą-* $\sigma(E',E)$,
gdzie $E$ utożsamiamy z obrazem kanonicznego zanurzenia
$J:E\to E''$, $J(x)(f)=f(x)$.
Wtedy $\sigma(E',E)\subseteq\sigma(E',E'')$,
a równość zachodzi wtedy i tylko wtedy, gdy $E$ jest
refleksywna (tj. $J$ jest surjekcją).
\end{stwierdzenie}
\begin{proof}
Pierwsza inkluzja wynika z $J(E)\subseteq E''$.
Jeżeli $J(E)=E''$, rodziny definiujących funkcjonałów
są jednakowe, więc topologie są równe.
Odwrotnie, jeżeli topologie są równe, ich przestrzenie
dualne topologiczne są równe. Z Twierdzenia
\ref{thm:slaby-dual}, zastosowanego do par
$(E',J(E))$ i $(E',E'')$, wynika odpowiednio,
że te przestrzenie dualne są algebraicznie równe
$J(E)$ i $E''$. Zatem $J(E)=E''$.
\end{proof}

\begin{przyklad}[Zbieżność słaba-* nie musi być słaba]
Niech $E=c_0$, $E'=\ell^1$ z naturalnym parowaniem.
Ciąg wektorów jednostkowych $e_n\in\ell^1$ spełnia
$e_n\to0$ w $\sigma(\ell^1,c_0)$, ponieważ dla każdego
$x\in c_0$ mamy $\langle x,e_n\rangle=x_n\to0$.
Nie jest to jednak zbieżność słaba w $\ell^1$:
funkcjonał z $\ell^\infty=(\ell^1)'$ odpowiadający
ciągowi $(1,1,\ldots)$ ma na każdym $e_n$ wartość $1$.
\end{przyklad}

\begin{przyklad}[Topologia produktowa jako topologia słaba]
Niech $E=\K^{\N}$ z topologią produktową i $F=c_{00}$
z parowaniem $\langle x,a\rangle=\sum_n x_na_n$
(suma jest skończona). Ponieważ współrzędne
$x\mapsto x_n$ odpowiadają wektorom $e_n\in F$,
topologia produktowa zawiera $\sigma(E,F)$ i odwrotnie:
każdy funkcjonał z $F$ jest skończoną kombinacją
odwzorowań współrzędnych. Zatem topologia produktowa
jest dokładnie $\sigma(E,F)$.
\end{przyklad}

\subsection{Zupełność topologii słabej i przestrzeń dualna algebraiczna}

\begin{twierdzenie}[Uzupełnienie przestrzeni ze słabą topologią]
\label{thm:slabe-uzupelnienie}
Niech $(E,F)$ będzie parą dualną. Niech $F^*$ oznacza
przestrzeń dualną algebraiczną przestrzeni liniowej $F$,
wyposażoną w topologię punktowej zbieżności na $F$.
Wówczas odwzorowanie $J:E\to F^*$, określone wzorem
$J(x)(f)=\langle x,f\rangle$, jest liniowym zanurzeniem
topologicznym przestrzeni $(E,\sigma(E,F))$ w $F^*$.
Obraz $J(E)$ jest gęsty, a $F^*$ jest zupełną TVS
Hausdorffa. Tym samym $F^*$ jest uzupełnieniem
$(E,\sigma(E,F))$ w sensie \hyperref[sec:zupelnosc]{rozdziału 3}.
\end{twierdzenie}
\begin{proof}
Traktujemy $F^*$ jako podzbiór produktu $\K^F$.
Jest to zbiór tych funkcji $\Phi:F\to\K$, które spełniają
$\Phi(f+g)=\Phi(f)+\Phi(g)$ i
$\Phi(af)=a\Phi(f)$ dla $f,g\in F$, $a\in\K$.
Każdy taki warunek definiuje zbiór domknięty w topologii
produktowej, zatem $F^*$ jest domkniętą podprzestrzenią
$\K^F$. Produkt zupełnych przestrzeni liniowo-topologicznych
jest zupełny, a podprzestrzeń domknięta przestrzeni
zupełnej jest zupełna. Topologia jest Hausdorffa.

Niezdegenerowaność parowania daje różnowartościowość $J$,
a definicje topologii początkowych oznaczają, że jest to
zanurzenie topologiczne. Aby wykazać gęstość, weźmy
$\Phi\in F^*$ i podstawowe otoczenie punktu $\Phi$,
wyznaczone przez $f_1,\ldots,f_m\in F$ oraz $\varepsilon>0$.
Niech $G=\operatorname{lin}\{f_1,\ldots,f_m\}\subseteq F$.
Odwzorowanie $T:E\to G^*$, $T(x)(f)=\langle x,f\rangle$,
jest surjekcją. Istotnie, gdyby jego obraz był właściwą
podprzestrzenią skończenie wymiarowej przestrzeni $G^*$,
istniałby niezerowy element $g\in G$ z
$\langle x,g\rangle=0$ dla każdego $x\in E$,
co przeczy niezdegenerowaności parowania.
Wybieramy zatem $x\in E$ takie, że
$T(x)=\Phi|_G$. Wtedy $J(x)(f_j)=\Phi(f_j)$
dla wszystkich $j$, a więc $J(x)$ należy do rozważanego
otoczenia. To dowodzi gęstości.
\end{proof}

\begin{wniosek}[Zaskakująca niezupełność topologii słabej]
\label{cor:slabe-niezupełne}
Jeśli $E$ jest nieskończenie wymiarową przestrzenią
unormowaną, to przestrzeń $(E,\sigma(E,E'))$ nie jest
zupełna, nawet gdy $E$ jest przestrzenią Hilberta.
\end{wniosek}
\begin{proof}
Przestrzeń $E'$ jest nieskończenie wymiarową przestrzenią
Banacha. Istnieje na niej nieciągły funkcjonał liniowy
$\Phi\in(E')^*$ (przykład z \hyperref[sec:hahn-banach]{rozdziału 6}).
Każdy $J(x)$, $x\in E$, jest jednak ciągły względem normy
operatorowej na $E'$, ponieważ
$|J(x)(f)|=|f(x)|\le\|x\|\|f\|$.
Stąd $\Phi\notin J(E)$. Twierdzenie
\ref{thm:slabe-uzupelnienie} mówi, że $J(E)$ jest
gęstą, właściwą podprzestrzenią zupełnej TVS Hausdorffa
$F^*=(E')^*$. Zatem $J(E)$ nie jest zupełne.
\end{proof}

\begin{uwaga}
Powyższy wniosek nie przeczy zupełności $E$ w normie:
zupełność zależy od topologii i jej struktury jednostajnej.
Nie należy też mylić przestrzeni dualnej algebraicznej
$(E')^*$ z bidualem topologicznym $E''=(E')'$.
\end{uwaga}

\subsection{Topologia silna i zapowiedź topologii Mackeya}

\begin{definicja}[Topologia silna na przestrzeni dualnej]
\label{def:silna-dualna}
Niech $E$ będzie lokalnie wypukłą TVS Hausdorffa.
Dla każdego niepustego ograniczonego zbioru $B\subseteq E$ i $f\in E'$
kładziemy $p_B(f)=\sup_{x\in B}|f(x)|$.
Każdy $p_B$ jest skończony, ponieważ ciągły funkcjonał
przekształca zbiory ograniczone na ograniczone zbiory skalarów.
Topologię na $E'$ generowaną przez rodzinę takich seminorm
nazywamy topologią silną i oznaczamy $\beta(E',E)$.
\end{definicja}
\begin{stwierdzenie}[Przypadek przestrzeni unormowanych]
\label{prop:silna-unormowana}
Jeśli $E$ jest przestrzenią unormowaną, to topologia
$\beta(E',E)$ jest topologią normy operatorowej
$\|f\|=\sup_{\|x\|\le1}|f(x)|$ na $E'$.
\end{stwierdzenie}
\begin{proof}
Dla ograniczonego $B$ mamy
$\sup_{x\in B}\|x\|=M<\infty$, więc
$p_B(f)\le M\|f\|$. Z drugiej strony kula jednostkowa
$B_E$ jest ograniczona, a $p_{B_E}(f)=\|f\|$.
Obie rodziny seminorm definiują zatem tę samą topologię.
\end{proof}

\begin{definicja}[Topologia Mackeya]\label{def:mackey}
Niech $(E,F)$ będzie parą dualną. Dla każdego absolutnie
wypukłego, niepustego i zwartego w $\sigma(F,E)$ zbioru $K\subseteq F$
określamy $p_K(x)=\sup_{f\in K}|\langle x,f\rangle|$.
Topologię lokalnie wypukłą generowaną przez wszystkie takie
seminormy nazywamy \emph{topologią Mackeya} $\tau(E,F)$.
\end{definicja}

\begin{stwierdzenie}[Poprawność definicji Mackeya]\label{prop:mackey-basic}
Dla dowolnej pary dualnej $(E,F)$ funkcje $p_K$ są
skończonymi seminormami, a topologia $\tau(E,F)$
jest lokalnie wypukła, Hausdorffa i silniejsza od
$\sigma(E,F)$.
\end{stwierdzenie}
\begin{proof}
Dla ustalonego $x\in E$ funkcja
$f\mapsto|\langle x,f\rangle|$ jest ciągła względem
$\sigma(F,E)$, a zatem ograniczona na każdym zwartym $K$.
Nierówność trójkąta i jednorodność dają własności seminormy.
Dla każdego $f\in F$ zbiór
$K_f=\{af:|a|\le1\}$ jest absolutnie wypukły i zwarty
jako obraz zwartego domkniętego koła w $\K$.
Mamy $p_{K_f}(x)=|\langle x,f\rangle|$.
Topologia Mackeya zawiera więc topologię słabą,
a ponieważ ta ostatnia jest Hausdorffa, taką własność
ma również $\tau(E,F)$.
\end{proof}

\begin{uwaga}[Twierdzenie Mackeya--Arensa: zapowiedź]
Twierdzenie Mackeya--Arensa orzeka, że $\tau(E,F)$ jest
najsilniejszą topologią lokalnie wypukłą na $E$, dla której
przestrzeń dualna topologiczna jest algebraicznie równa $F$
(przy podanym parowaniu). Wymaga to jednak twierdzeń
o zwartości i zbiorach polarnych, które rozwiniemy w następnym
rozdziale. Nie traktujemy tej charakterystyki jako
udowodnionej w bieżącym miejscu.
\end{uwaga}

\begin{uwaga}[W stronę dystrybucji]
Dla przestrzeni funkcji testowych $\mathcal D(\Omega)$
topologia słaba-* $\sigma(\mathcal D'(\Omega),\mathcal D(\Omega))$
opisuje zbieżność dystrybucji na każdej ustalonej funkcji
testowej. Topologia silna przestrzeni dualnej opisuje
natomiast zbieżność jednostajną na wszystkich zbiorach
ograniczonych funkcji testowych. Ich porównanie będzie
istotne w dalszych rozdziałach.
\end{uwaga}
